\subsection{Modular obstructions to a three-plane}
\label{sec:gct-modular}

Let \(V_R\) be the symmetric entry-sum-zero tangent space over
\(R=\mathbb Z_{(p)}\), for a prime \(p>3\). The level locus for
\(r\)-planes in \(V_R\) is closed in \(\operatorname{Gr}(r,V_R)\).
Its equations are the restrictions of \(q_2,q_3,q_4\) to the
tautological subbundle. The same construction applies inside a
coordinate subspace of \(V_R\) in Hadamard coordinates.
The incidence scheme is projective over \(R\), hence proper
\cite{stacksprojectiveproper}. If its geometric special fiber is
empty, its closed image in \(\operatorname{Spec}R\) misses the closed
point and therefore also the generic point. This excludes complex
points. When only the set of \(\mathbb F_p\)-points is empty, the
valuative criterion excludes points over any local field of
characteristic zero with residue field \(\mathbb F_p\)
\cite{stacksvaluativeproper}.

\paragraph{Rational three-planes.}
For \(p=5,7\), keep the three rational Witt pairs from
Section~\ref{sec:gct-three-plane}. Choose \(a,b\in\mathbb F_p\) with
\(1+a^2+b^2=0\). In the span of the three orthogonal vectors \(w_i\),
take the fourth pair to be
\(w_1+a w_2+b w_3\), \((w_1-a w_2-b w_3)/24\).
The remaining vector \(-b w_2+a w_3\) has norm \(-12\).
The same projection-minor argument gives all 32 charts.
In a normalized frame, each column has three free complement coordinates
and two free dual coordinates. Its own dual coordinate is fixed by
\(q_2=0\). Enumerate these five entries over \(\mathbb F_p\) for each
column and retain only columns with \(q_3=q_4=0\). Join choices in
distinct columns when their mixed quadratic coefficient vanishes and
when \(q_3,q_4\) vanish on their sums with relative weights \(1,2,3\).
Every level three-plane would give a triangle in this graph.
The executable finds no triangles in either prime field.

Thus the level locus has no \(\mathbb Q_5\)- or \(\mathbb Q_7\)-point,
and in particular no rational three-plane. Any number field of definition
must have no residue-degree-one prime above either 5 or 7. This also
excludes \(\mathbb Q(\mathrm i)\) as a field of definition, since 5 splits
there. The computation concerns points over the prime fields.

\paragraph{Rank of the row-sum map.}
A three-plane on which \(B\mapsto B\mathbf1\) has rank at most one
contains a two-plane in the zero-row-sum slice. To exclude such
two-planes, work over \(\overline{\mathbb F}_{11}\), where the quadratic
form on this slice is \(\sum_i(x_i^2+2y_i^2)\) and \(\kappa=4\)
satisfies \(\kappa^2=-1/2\). Of the three hyperbolic pairs in
coordinates \((x_i,y_i)\), choose two pivot lines. Simultaneous
permutations of the three Klein characters permute the pairs.
A Klein translation interchanges either zero or two pairs of lines
and can swap just one of the selected pairs by also swapping the
omitted pair. So one normalized chart covers the existence question.
With frame parameters \(s,t\), put
\begin{align*}
 u&=as+bt,& v&=cs+dt,\\
 \alpha&=gt-(au-cv)/4,& \beta&=-gs-(bu-dv)/4,\\
 x&=(s+\alpha,t+\beta,u),& y&=\kappa(s-\alpha,t-\beta,v).
\end{align*}
The pivot coordinates are \((x_1+y_1/\kappa)/2=s\) and
\((x_2+y_2/\kappa)/2=t\). The symmetric part of the two dual
coordinates is forced by \(4(s\alpha+t\beta)+u^2-v^2=0\), leaving the
one skew parameter \(g\).
The four cubic coefficients of \(\operatorname{per}(C)\) and five
quartic coefficients of \(\operatorname{per}(HCH^T)\) generate the unit
ideal over \(\mathbb F_{11}\). The checker reaches a nonzero constant by
modular row reduction and multiplication by variables within total degree
ten. Every row is in the ideal generated by those nine coefficients.
This proves geometric emptiness of the special fiber. Consequently every
complex level three-plane has row-sum map of rank at least two.

\paragraph{Rank of the Hadamard diagonal map.}
In the coordinates of Section~\ref{sec:gct-three-plane}, consider
\(B\mapsto(d_1,d_2,d_3)\), the diagonal of the lower Hadamard block.
A rank at most one on a level three-plane would give a level two-plane
with \(d=0\). On that slice put \(u_i=c_i+b_i\) and \(v_i=c_i-b_i\).
The permanent coefficients are
\begin{align*}
 q_2/2 &= \sum_i(c_i^2-3b_i^2),&
 q_3/2 &= \prod_i v_i+4\prod_i b_i,\\
 16q_4 &= \left(\sum_i v_i^2-\sum_i u_i^2\right)^2
          +4\sum_i v_i^2u_i^2-4\sum_{i<j}u_i^2u_j^2.
\end{align*}
The checker verifies the identities by exact expansion.

Over \(\overline{\mathbb F}_{23}\), first suppose every \(b_i\) is a
nonzero linear form on the two-plane. Unique factorization applied to
\(q_3=0\) yields \(v_i=k_i b_{\sigma(i)}\) with \(\prod_i k_i=-4\).
Then \(b\) determines all the coordinates, so its rank on the plane is
two. Relabeling the characters gives \(b=(s,t,\alpha s+\beta t)\).
Swapping the first two labels leaves four permutation representatives
\(1,(12),(13),(123)\). The checker obtains the unit ideal on each
branch from the product, quadratic, and quartic equations.

Otherwise one can take \(b_1=0\), and \(q_3=0\) forces some
\(c_j=b_j\). Up to relabeling the other indices, \(j=1\) or \(j=2\).
The restricted quadratic form \(q_2/2\) is then
\(c_2^2-3b_2^2+c_3^2-3b_3^2\) or
\(c_1^2-2b_2^2+c_3^2-3b_3^2\).
Each is the sum of two forms \(z^2-\lambda w^2\) with
\(\lambda\in\{2,3\}\), both squares modulo 23. Each of the four
projection charts is parametrized by one skew entry. The quartic
coefficients have greatest common divisor one in all eight charts.
The two branches cover the geometric special fiber, so every complex
level three-plane has Hadamard diagonal map of rank at least two.
The full existence problem over \(\mathbb C\) remains open.
